\chapter{Tests}
\label{chap:Tests}

\begin{tcolorbox}[colback=aliceblue!10!white,colframe=blue!70!black,title=\textbf{Remark}]
% olback=green!10!white,colframe=green!80!black
There is a tendency to use hypothesis testing methods even
when they are not appropriate. Often, estimation and confidence intervals are
better tools. Use hypothesis testing only when you want to test a well-defined
hypothesis \citep{wasserman2004all}.
\end{tcolorbox}


\begin{tcolorbox}[colback=aliceblue!10!white,colframe=green!90!black,title=\textbf{Remark}]
Remember that statistical methods cannot prove that a factor (or factors) has a particular effect. They only provide guidelines as to the reliability and validity of results.
Properly applied, statistical methods do not allow anything to be proved experimentally,
but they do allow us to measure the likely error in a conclusion or to attach a level of
confidence to a statement. The primary advantage of statistical methods is that they add
objectivity to the decision-making process. Statistical techniques coupled with good 
engineering or process knowledge and common sense will usually lead to sound
conclusions \citep{montgomery2017design}.
\end{tcolorbox}

\section{Introduction}
\label{sec:TestsIntroduction}

In this chapter we will cover tests such as t-test and A/B testing.



\section{A/B Testing}
\label{sec:ABTesting}

Like anything and everything in statistics there is no unique definition
for a unique idea. A/B testing in some fields might be called controlled experiment,
or randomized controlled experiment and so on.
The concept is simple; we have an idea and we want to know
if this idea makes a difference compared to the status quo
and if so, by how much. 
In the field of medicine it will be testing a new 
drug and in online setting it would be testing a new feature on the 
app, etc. We divide our \emph{units} into two groups; one group
is called \emph{control} which sees no change (i.e. it will not receive any update/drug/new feature),
and there is another group called \emph{treatment} which will see the new change (get the new drug,
new feature).


In order to do so, we need to divide our units (could be people, animals, codes, computers, etc.)
into two groups. 
These memberships are assigned randomly unless there are other considerations.
For example, we divide online users into 2 groups randomly
and this randomization is not tied to any other things such as demographics,
or membership of one person to control group will not affect membership of another
person to any of these groups.

But, before proceeding to the A/B testing, an A/A testing is done.
In A/A testing none of the groups receive the treatment. They both will
see/receive the same features/drugs.
The purpose of A/A testing is to investigate the set up of the experiment;
\begin{itemize}
\setlength\itemsep{0pt}

\item is the randomization done correctly? Is the population divided into 50-50 groups
or there are 52\% in one group and 48\% in another. This 2\% can be a huge difference
and should be prevented according to Kohavi in a podcast, and I am sure it will come up
in his book\cite{kohavi2020trustworthy} (which I just started to read currently).

\item Is the randomization done with no bias?

\item Are the users' activity from each group being logged correctly?

\item Establish a baseline (observe natural variation before any treatment)

\end{itemize}


Let me step back even further. Before any experimental design,
there should be an objective that is well-defined, measurable, and agreed upon
by parties involved\sidenote{It happens that in some setting more than one party might be involved and in another setting perhaps there is only one experimenter.}. 
It turns out a new feature on an app
can be tied to other systems such as security, user-experience, etc.
For example, the whole page can be filled with ads, and revenue go up (at least in short term),
but user-experience will be affected negatively. 
Other than interconnected issues, there is a problem of short-term versus
long term results/thinkiong (which is usually ignored in real life).
So, the chosen metric must be well-established and thought of.
Even for short-term results a good metric must be chosen.
For example, click through rate can be high, but after clicking,
users can just close the new window, and user-session (engagement length) can be a
better metric perhaps.
According to Kohavi metrics for tests should be measurable in short term (during the 2(?) weeks experiment)
and be a good representative for long-term goals.
For example, if click through rate is the metric, perhaps a video that goes
viral will be recommended, but after a few weeks, click through rate will drop.
(Or in case of movie streaming platforms, you cannot recommend the same thing
again and again. A balance between exploration and exploitation must be considered.)\\

Let me go back to A/B testing.
Another purpose of A/B testing is to establish causality.
Observational studies can establish only correlation.
This is said repeatedly everywhere. Let me mention an
example from\cite{kohavi2020trustworthy}. 
There were data from Office 365. Users with lots of crashes and
bug reports had more engagement with the App. This does not mean
Microsoft should try to increase crashes to get more customers.
It just happens that users who use Office 365 more often, get more crashes.

Other than metrics, and randomization of users, we must have access to large
population. Especially if we want to detect small differences.

\begin{tcolorbox}[colback=aliceblue!10!white,colframe=blue!70!black,title=\textbf{A short break}]
In A/B testing one idea is tried at a time. At least, that is my understanding
at this moment. In tech companies there are thousands of ideas to try
and it seems they do it one idea at a time.
But, experimental design can have different purposes with
more than one variable involved. 
One variable at a time, fails to consider interaction of different variables 
with each other.
If there are two ideas each of which with two levels, their combinations will have 4 distinct 
levels and consequently 4 experiments are needed. 
Factorial Design is like grid search. From Page 7 of Montgomery's 5th Edition book: 
``if there are four to five or more factors, it is usually unnecessary to run
all possible combinations of factor levels. A fractional factorial experiment is a variation of the basic factorial design in which only a subset of the runs are made.''
Side note about experimentation: from \emph{Design and Analysis of Experiments}
by Montgomery on designing experiments (not A/B testing):
\begin{itemize}
\setlength\itemsep{0pt}
\item Determine which variables are most influential
\item Where to set independent variables so that the output is at/near the desired target
\item What should input variables be so that the output has minimal variance
\item Where should controllable variables be set so that influence of uncontrollable variables are minimized
\end{itemize}
\end{tcolorbox}




\begin{tcolorbox}[colback=aliceblue!10!white,colframe=blue!70!black,title=\textbf{Follow the Money}]
It's all about money and control. 
The purpose of A/B and lots of experiments is about money
and \emph{optimization}.
\begin{itemize}
\setlength\itemsep{0pt}
\item improve yield
\item improve cost
\item improve development time
\item reduce variability of response
\end{itemize}
\end{tcolorbox}

